\subsection{MAS-10: an ordinary Xi determinant on the actual weighted quotient}
The trace-ideal question for the reciprocal-square multiplier has an affirmative answer. The extra ingredient is a sector estimate in the actual weighted Hilbert norm. It avoids any assumption that the zero modes form a Riesz basis or have uniformly bounded biorthogonal vectors.

Specialize to $F=\mathbf Q$, the trivial character of $C_{\mathbf Q}^1$, center $a=1/2$, width $\alpha=1/2$, and $s>1$. Normalize the compact Haar factor to mass one and write $u=\log|g|$. Thus
\[
 H=L^2(\R,w(u)\,du),\qquad
 w(u)=e^{|u|}(1+u^2)^{s/2},\qquad E=H/V,
\]
where $V$ is the closed centered theta-image. This is the actual Hilbert quotient of \cite[(10.2.2.1)--(10.2.2.2), p.~445]{masullo}, in the indicated character sector.

\textbf{Exact conditional input.} Use \cite[Theorem 10.2.2.1, (10.2.2.4)--(10.2.2.6), p.~446]{masullo} only for its full spectral-annihilator description. Let $\rho$ range over distinct nontrivial zeta zeros, with multiplicity $N_\rho$, and put $z_\rho=\rho-1/2$. Classically
\[
 0<\Re\rho<1,\qquad z_\rho\ne0,
\]
by the Euler product, the functional equation, nonvanishing on $\Re s=1$, and $\xi(1/2)\ne0$; see \cite[\S2.12; \S3.2, pp.~46--47]{titch}. These facts require no RH. Every zero is therefore interior to the chosen strip, so there is no boundary truncation. The cited draft claim says that
\[
 \ell_{\rho,m}(f)=\int_{\R}f(u)e^{z_\rho u}u^m\,du,
 \qquad 0\le m<N_\rho,
\]
annihilate $V$ and have norm-dense linear span in the Hilbert dual $E'$. This Hilbert-dual density is also explained in \cite[Remark 10.2.1.2, pp.~440--441]{masullo}. No nuclearity assertion from Theorem 10.3.2.1 is used. The conclusions below are conditional on the spectral-annihilator claim alone, together with the stated classical zeta facts; \emph{RH is not an input}.

Put $B_0=1/2+s/4$ and choose an auxiliary real parameter $\mu>B_0$, distinct from the center $a=1/2$. On the ambient space define
\[
 g_\mu(v)=\frac{e^{-\mu|v|}}{2\mu},\qquad
 K_\mu f=\int_{\R}g_\mu(v)f(\,\cdot-v)\,dv.
\]
The integral is a strong, vector-valued integral for each $f$, as justified below. Let $K$ be its induced operator on $E$.

\begin{proposition}
The operator $K$ is trace class and $I-\mu^2K$ has a bounded inverse. The operator
\[
 R=K(I-\mu^2K)^{-1}
\]
is trace class, is independent of $\mu$, and has the full jet multiplier $-1/z^2$ at $z=z_\rho$. Reflection $Jf(u)=f(-u)$ descends to a unitary involution of $E$ and commutes with $R$. For $E_+=\ker(J-I)$ and $R_+=R|_{E_+}$,
\begin{equation}\label{eq:masactualXi}
 \det_{E_+}(I-wR_+)=\frac{\Xi(\sqrt w)}{\Xi(0)}.
\end{equation}
Each pair $\{\rho,1-\rho\}$ contributes its full multiplicity $N_\rho$, exactly once. Moreover
\begin{equation}\label{eq:masGreennorm}
 \|K\|_1\le\frac{\mu}{\sqrt{\mu^2-B_0^2}}
 \sum_\rho N_\rho\Re\frac1{\mu^2-z_\rho^2},
\end{equation}
where the sum includes all distinct nontrivial zeros, and every real part in it is positive.
\end{proposition}
\begin{proof}
\emph{Translations and the quotient.}
The function $\log\sqrt w$ is Lipschitz, with almost-everywhere derivative
\[
 b(u)=\tfrac12\operatorname{sgn}(u)+\frac{s u}{2(1+u^2)},
 \qquad |b(u)|\le B_0.
\]
Consequently translations $T_vf(u)=f(u-v)$ satisfy
$\|T_v\|\le e^{B_0|v|}$. They are strongly continuous: first check this on compactly supported smooth functions, and then use their density and the locally uniform bound. The functionals $\ell_{\rho,m}$ are continuous because
$\int_{\R}e^{2\Re(z_\rho)u}|u|^{2m}w(u)^{-1}\,du<\infty$. Substitution gives
\begin{equation}\label{eq:masjettranslation}
 \ell_{\rho,m}(T_vf)=e^{z_\rho v}
 \sum_{j=0}^m\binom{m}{j}v^{m-j}\ell_{\rho,j}(f).
\end{equation}
The complete annihilator description therefore implies $T_vV\subset V$: every jet vanishes on $T_vf$ when $f\in V$, and $V$ is the annihilator of its annihilator. Since
$g_\mu\in L^1(\R,e^{B_0|v|}\,dv)$, the strong integral defining $K_\mu$ is bounded and preserves $V$. Under the isometric identification $E=V^\perp$, its quotient operator is the compression
\[
 K=P_{V^\perp}K_\mu|_{V^\perp}.
\]
No operator-norm Bochner integrability of $v\mapsto T_v$ is asserted or required.

\emph{The weak domain and the sector estimate.}
Both $g_\mu$ and $g_\mu'$ belong to the same exponentially weighted $L^1$ space, and
$g_\mu''=\mu^2g_\mu-\delta_0$ in distributions. Hence $y=K_\mu f$ satisfies
\[
 y,y',y''\in H,\qquad y''=\mu^2y-f.
\]
These statements follow first for compactly supported smooth $f$ by convolution, then for arbitrary $f\in H$ by the weighted convolution bounds and convergence in distributions. Let $U:H\longrightarrow L^2(\R)$ be multiplication by $\sqrt w$ and put $h=Uy$. Then
\[
 h\in H^1(\R),\quad D h:=h'-bh=Uy'\in H^1(\R),
 \quad D^2h=Uy''.
\]
Indeed $h'=\sqrt w\,y'+bh$, and the same calculation applies to $y'$. The weak derivative of $b$ is never used, so the cusp of the weight at zero causes no omitted boundary term. Integration by parts in $H^1$ gives
\begin{align*}
 \langle(\mu^2-D^2)h,h\rangle
 &=\|h'\|_2^2+\mu^2\|h\|_2^2-\|bh\|_2^2\notag\\
 &\quad+\langle h',bh\rangle-\langle bh,h'\rangle.
\end{align*}
For example, $D^*=-\partial_u-b$ on $H^1$, and the identity follows from
$-\langle Dh,D^*h\rangle$; this also justifies it on the stated weak domain. With
$\delta=\mu^2-B_0^2>0$ and $c=B_0/\sqrt\delta$, it follows that
\[
 \Re\langle(\mu^2-D^2)h,h\rangle
 \ge\|h'\|_2^2+\delta\|h\|_2^2,
\]
\[
 |\Im\langle(\mu^2-D^2)h,h\rangle|
 \le2B_0\|h'\|_2\|h\|_2
 \le c\bigl(\|h'\|_2^2+\delta\|h\|_2^2\bigr).
\]
The last inequality is $2XY\le X^2+Y^2$. Since
$(\mu^2-d^2/du^2)y=f$, the number $\langle K_\mu f,f\rangle$ is the conjugate of
$\langle(\mu^2-d^2/du^2)y,y\rangle$. Thus
\begin{equation}\label{eq:masGreensector}
 \Re\langle K_\mu f,f\rangle\ge0,
 \qquad |\Im\langle K_\mu f,f\rangle|
 \le c\Re\langle K_\mu f,f\rangle.
\end{equation}
Compression preserves this numerical-range inequality because
$\langle PK_\mu x,x\rangle=\langle K_\mu x,x\rangle$ for $x\in V^\perp$.
Therefore both $K$ and $K^*$ satisfy the same symmetric sector bound.

\emph{A trace-ideal lemma.}
If a bounded operator $T$ satisfies the sector bound in \eqref{eq:masGreensector} and an orthonormal basis $(e_n)$ satisfies
$\sum_n\Re\langle Te_n,e_n\rangle<\infty$, then $T$ is trace class. To prove this, set
\[
 A=(T+T^*)/2\ge0,\qquad Q=(T-T^*)/(2i).
\]
The diagonal sum says
$\sum_n\|A^{1/2}e_n\|^2<\infty$, so $A^{1/2}$ is Hilbert--Schmidt and $A$ is trace class. The sector bound says $-cA\le Q\le cA$ in quadratic forms. It yields
\[
 Q=A^{1/2}CA^{1/2}
\]
for a bounded selfadjoint $C$ with $\|C\|\le c$. Here is the factorization argument. Define on $\operatorname{Ran}A^{1/2}$ the form
$\beta(A^{1/2}x,A^{1/2}y)=\langle Qx,y\rangle$.
It is well defined: if $x\in\ker A$, the positive operators $cA\pm Q$ have zero quadratic value at $x$ and therefore annihilate $x$, implying $Qx=0$. Cauchy--Schwarz for these two positive forms makes $\beta$ bounded. Its diagonal bound
$|\beta(u,u)|\le c\|u\|^2$ extends to the closure and represents a bounded selfadjoint operator $C$ of norm at most $c$. Extend it by zero on the orthogonal complement. Consequently
\[
 T=A^{1/2}(I+iC)A^{1/2},\qquad
 \|T\|_1\le\|I+iC\|\|A^{1/2}\|_2^2
 \le\sqrt{1+c^2}\,\Tr A.
\]
The trace-class estimate follows either from the Hilbert--Schmidt product inequality or directly by expanding the product in an orthonormal basis and applying Cauchy--Schwarz to the sum of the rank-one norms. This proves the lemma without a prior spectral trace identity.

\emph{Full jets and an orthonormal triangular basis.}
For $|\Re z|<\mu$,
\[
 F_\mu(z):=\int_{\R}g_\mu(v)e^{zv}\,dv
 =\frac1{\mu^2-z^2}.
\]
Integrating \eqref{eq:masjettranslation}, with absolutely convergent exponential-polynomial moments, gives
\[
 K'\ell_{\rho,m}
 =\sum_{j=0}^m\binom{m}{j}
 F_\mu^{(m-j)}(z_\rho)\ell_{\rho,j},
 \qquad \lambda_\rho:=F_\mu(z_\rho)=\frac1{\mu^2-z_\rho^2}.
\]
These eigenvalues need not be real. If $z_\rho=b+it$, set $A_\rho=\mu^2-b^2+t^2>0$. Then
\[
 \Re\lambda_\rho=\frac{A_\rho}{A_\rho^2+4b^2t^2}>0,
 \qquad |\lambda_\rho|\le\frac1{\mu^2-1/4+t^2}.
\]
The classical zero count used in \fid{B-01} consequently gives
$\sum_\rho N_\rho|\lambda_\rho|<\infty$.

Enumerate the distinct zeros and list each finite jet block in increasing jet order. Their Riesz representatives form a total linearly independent sequence in $E$, and every finite prefix spans a $K^*$-invariant space. Applying Gram--Schmidt preserves these finite-prefix spaces and gives an orthonormal basis $(e_n)$ such that, in the paper's conjugate-linear-first convention,
\[
 K^*e_n-\overline{\lambda_n}e_n\in\operatorname{span}(e_1,\ldots,e_{n-1}),
 \qquad \langle e_n,K^*e_n\rangle=\overline{\lambda_n}.
\]
The conjugation is required by the anti-linear Riesz identification. Each $\lambda_\rho$ occurs $N_\rho$ times. No bounds on the original basis vectors or their biorthogonal partners enter this argument. Hence
\[
 \sum_n\Re\langle K^*e_n,e_n\rangle
 =\sum_\rho N_\rho\Re\lambda_\rho<\infty.
\]
The trace-ideal lemma and the sector bound prove $K\in\Sone(E)$ and
\eqref{eq:masGreennorm}. This establishes trace class before using compactness or a trace-equals-eigenvalue theorem; summable eigenvalues of a general nonnormal operator have not been substituted for singular-value control.

\emph{Exact spectrum, inverse and the reciprocal-square multiplier.}
Since $K^*$ is now compact, if $P_n$ projects onto the first $n$ basis vectors, then
$\|(I-P_n)K^*\|\to0$. To see this, approximate $K^*$ in norm by a finite-rank operator and use uniform convergence of $P_n$ on its finite-dimensional range. The invariant flags give a block matrix
\[
 K^*=\begin{pmatrix}T_n&B_n\\0&D_n\end{pmatrix},
 \qquad \|D_n\|\longrightarrow0,
\]
with $T_n$ finite triangular of diagonal $\overline{\lambda_1},\ldots,\overline{\lambda_n}$.
If $q\ne0$ differs from all these diagonal values, choose $n$ with $\|D_n\|<|q|$. Both $T_n-qI$ and $D_n-qI$ are invertible, the latter by a Neumann series, and the explicit upper-triangular block inverse makes $K^*-qI$ invertible. Conversely each diagonal value is an eigenvalue in a finite invariant prefix. For a fixed nonzero eigenvalue, take $n$ past all its occurrences with $\|D_n\|$ less than its modulus. Every generalized eigenvector then lies in that prefix, since the tail block minus the eigenvalue is invertible. Its algebraic multiplicity is exactly the finite triangular multiplicity. There are only finitely many occurrences: equality $\lambda_\rho=\lambda_{\rho'}$ is equivalent to $z_\rho^2=z_{\rho'}^2$, hence to $\rho'=\rho$ or $1-\rho$. Thus the nonzero spectrum of $K$, with algebraic multiplicities, is precisely the $\lambda_\rho$; no additional nonzero spectral part remains.

Since $z_\rho\ne0$, no $\lambda_\rho$ equals $1/\mu^2$, so $I-\mu^2K$ is boundedly invertible. The ideal property now proves
$R=K(I-\mu^2K)^{-1}\in\Sone(E)$. On every finite jet block, rational composition gives the multiplier
\[
 \frac{F_\mu(z)}{1-\mu^2F_\mu(z)}=-\frac1{z^2}.
\]
The identity holds with all jets at $z=z_\rho$, since the denominator is nonzero there; equivalently it is an identity in the corresponding truncated Taylor algebra. Its diagonal is $-z_\rho^{-2}$ with full multiplicity. The same compact triangular argument gives exactly these nonzero spectral values and their algebraic multiplicities. Two choices of $\mu$ act identically on every dual jet, so their difference is annihilated by a dense subspace of $E'$ and is zero. This proves independence of the auxiliary parameter.

\emph{Parity and the determinant.}
Reflection $Jf(u)=f(-u)$ is a unitary involution on $H$, since $w$ is even, and
\[
 \ell_{\rho,m}(Jf)=(-1)^m\ell_{1-\rho,m}(f).
\]
The functional equation gives $N_{1-\rho}=N_\rho$ and $z_{1-\rho}=-z_\rho$. The complete annihilator description therefore yields $JV=V$, so $J$ descends to a unitary involution of $E$. The even Green kernel commutes with reflection, as do $K$ and $R$. Hence $R_+=R|_{E_+}$ is trace class on the orthogonal even summand.

For each pair $\{\rho,1-\rho\}$, reflection exchanges two independent jet spaces of dimension $N_\rho$ with the indicated signs. Its $+1$ eigenspace has dimension $N_\rho$. Applying $(I+J')/2$ to the dense full jet span proves density in the even dual. Thus $R_+$ has precisely one eigenvalue $-z_\rho^{-2}$ per pair, with algebraic multiplicity $N_\rho$. If $\gamma_\rho=-iz_\rho$, then $\gamma_\rho$ is a possibly complex zero of $\Xi$ and $\gamma_\rho^{-2}=-z_\rho^{-2}$. The ordinary trace-class determinant theorem of cote~1 and \fid{B-01} give
\[
 \det_{E_+}(I-wR_+)
 =\prod_{\{\rho,1-\rho\}}
 \left(1+\frac{w}{z_\rho^2}\right)^{N_\rho}
 =\frac{\Xi(\sqrt w)}{\Xi(0)}.
\]
The product is absolutely locally uniformly convergent. This proves every assertion without RH.
\end{proof}

The construction in \fid{MAS-10} is an ordinary Fredholm determinant on the stated arithmetic theta quotient, conditional on its full spectral-annihilator description. It does not assert that $R_+$ is selfadjoint or positive in that quotient metric, and equality of its determinant with the determinant of $B$ does not supply a unitary intertwiner. If, under RH, the heat construction of \fid{MAS-08} is made on this same quotient, its reciprocal-square operator equals $R$ because its full action on the separating dual jets is the same. No unresolved small-time trace-norm condition remains for the ordinary determinant in \eqref{eq:masactualXi}.
